\chapter{Parazwartość i podziały jedności}
\label{ch:topog-parazwartosc}

\noindent\textbf{Cel rozdziału.}
Własność Lindelöfa pozwala wybrać przeliczalne podpokrycie,
ale nie kontroluje liczby zbiorów spotykających małe otoczenie
danego punktu. Parazwartość wymaga lokalnie skończonych
uszczegółowień pokryć. Powiążemy ją z podziałami jedności
i skonstruujemy takie podziały dla regularnych przestrzeni
Lindelöfa, w tym dla ważnych klas przestrzeni metryzowalnych
oraz lokalnie zwartych. Korzystamy wyłącznie z wyników
rozdziałów~5, 9, 10 i 14.

\section{Pokrycia lokalnie skończone}
\label{sec:topog-parazwartosc-pokrycia}

\begin{definicja}[Rodzina lokalnie skończona]
\label{def:topog-parazwartosc-lokalnie-skonczona}
Niech $X$ będzie przestrzenią topologiczną, a
$(A_i)_{i\in I}$ rodziną jej podzbiorów.
Rodzina jest \emph{lokalnie skończona}, jeżeli
dla każdego $x\in X$ istnieje otwarte otoczenie
$O_x$ takie, że zbiór
$\{i\in I:O_x\cap A_i\ne\varnothing\}$
jest skończony. Powtórzenia indeksowanych
zbiorów są liczone osobno.
\end{definicja}

\begin{definicja}[Uszczegółowienie i parazwartość]
\label{def:topog-parazwartosc-definicja}
Niech $\mathcal U$ będzie pokryciem otwartym
przestrzeni $X$. Pokrycie $\mathcal V$ jest
\emph{uszczegółowieniem} $\mathcal U$, gdy
każdy $V\in\mathcal V$ zawiera się w pewnym
$U\in\mathcal U$. Przestrzeń $X$ nazywamy
\emph{parazwartą}, gdy każde jej pokrycie
otwarte ma lokalnie skończone uszczegółowienie
otwarte.
Nie zakładamy tu automatycznie
aksjomatu Hausdorffa; w twierdzeniach
wymagających go podajemy go jawnie.
\end{definicja}

\begin{stwierdzenie}[Domknięcia rodzin lokalnie skończonych]
\label{prop:topog-parazwartosc-domkniecia}
Niech $(A_i)_{i\in I}$ będzie lokalnie skończoną
rodziną podzbiorów przestrzeni topologicznej $X$.
Wówczas $(\overline{A_i}^{\,X})_{i\in I}$
jest lokalnie skończona, a ponadto
\[
\overline{\bigcup_{i\in I} A_i}^{\,X}
 =\bigcup_{i\in I}\overline{A_i}^{\,X}.
\]
\end{stwierdzenie}
\begin{proof}
Niech $x\in X$ i niech otwarty $O\ni x$
przecina tylko zbiory $A_i$ o indeksach
należących do skończonego $F\subseteq I$.
Jeżeli $i\notin F$, to $O\cap A_i=\varnothing$;
ponieważ $O$ jest otwarty, również
$O\cap\overline{A_i}^{\,X}=\varnothing$.
Zatem rodzina domknięć jest lokalnie skończona.

Inkluzja od prawej do lewej w podanej
równości wynika z monotoniczności domknięcia.
Dla odwrotnej ustalmy
$x\notin\bigcup_i\overline{A_i}^{\,X}$.
Wybierzmy otwarty $O\ni x$ przecinający
jedynie $A_i$, $i\in F$, gdzie $F$ jest
skończony. Zbiór
$O'=O\setminus\bigcup_{i\in F}
\overline{A_i}^{\,X}$
jest otwartym otoczeniem $x$ i nie
przecina żadnego $A_i$.
Stąd $x\notin\overline{\bigcup_i A_i}^{\,X}$.
\end{proof}

\begin{przyklad}[Lokalna skończoność nie oznacza skończoności]
\label{ex:topog-parazwartosc-przedzialy}
W $\R$ rodzina
$A_n=(n-\frac13,n+\frac13)$,
$n\in\mathbb Z$, jest nieskończona
i lokalnie skończona.
Istotnie, dla $x\in\R$ wybieramy
$\varepsilon\in(0,1)$.
Jeśli $(x-\varepsilon,x+\varepsilon)$
spotyka $A_n$, to
$|n-x|<\varepsilon+\frac13$,
co zachodzi tylko dla skończenie
wielu liczb całkowitych $n$.
W przeciwieństwie do tego rodzina
$(-1/n,1/n)$, $n\ge1$, nie jest
lokalnie skończona w punkcie $0$:
każde otoczenie $0$ przecina
wszystkie jej elementy.
\end{przyklad}

\section{Podziały jedności}
\label{sec:topog-parazwartosc-jednosc}

\begin{definicja}[Nośnik i podział jedności]
\label{def:topog-parazwartosc-podzial}
Niech $X$ będzie przestrzenią topologiczną.
Nośnikiem funkcji $\varphi:X\to\R$ nazywamy
$\operatorname{supp}\varphi
=\overline{\{x\in X:\varphi(x)\ne0\}}^{\,X}$.
\emph{Podziałem jedności} na $X$ nazywamy
rodzinę ciągłych funkcji
$\varphi_i:X\to[0,1]$, $i\in I$, dla której:
\begin{enumerate}[label=(\roman*)]
\item $(\operatorname{supp}\varphi_i)_{i\in I}$
      jest lokalnie skończona;
\item $\sum_{i\in I}\varphi_i(x)=1$
      dla każdego $x\in X$.
\end{enumerate}
Suma jest punktowo skończona na mocy (i).
Podział jest \emph{podporządkowany}
pokryciu otwartemu $\mathcal U$,
jeśli dla każdego $i$ istnieje
$U_i\in\mathcal U$ z
$\operatorname{supp}\varphi_i\subseteq U_i$.
\end{definicja}

\begin{stwierdzenie}[Podział jedności daje uszczegółowienie]
\label{prop:topog-parazwartosc-jednosc-pokrycie}
Jeżeli każde pokrycie otwarte przestrzeni
topologicznej $X$ ma podporządkowany
mu podział jedności, to $X$ jest parazwarta.
\end{stwierdzenie}
\begin{proof}
Niech $\mathcal U$ będzie pokryciem otwartym
i $(\varphi_i)$ podporządkowanym mu
podziałem jedności. Zbiory
$V_i=\{x:\varphi_i(x)>0\}$ są otwarte
na mocy ciągłości. Pokrywają $X$,
ponieważ $\sum_i\varphi_i(x)=1$
wyklucza jednoczesne zerowanie
wszystkich składników w punkcie $x$.
Każdy $V_i$ zawiera się w nośniku
$\varphi_i$, a przez to w pewnym
elemencie $\mathcal U$.
Rodzina $(V_i)$ jest lokalnie skończona,
gdyż każda jej część zawiera się
w odpowiadającym nośniku.
\end{proof}

\begin{lemat}[Funkcje odcinające w przestrzeni normalnej]
\label{lem:topog-parazwartosc-odciecia}
Niech $X$ będzie normalną przestrzenią
Hausdorffa, a $F\subseteq V\subseteq X$,
gdzie $F$ jest domknięty, a $V$ otwarty.
Istnieje funkcja ciągła $f:X\to[0,1]$
taka, że $f|_F=1$ oraz
$\operatorname{supp}f\subseteq V$,
o ile istnieje otwarty $W$ spełniający
$F\subseteq W\subseteq\overline W^{\,X}
\subseteq V$.
\end{lemat}
\begin{proof}
Zbiory $F$ i $X\setminus W$ są
rozłączne i domknięte. Lemat Urysohna
z rozdziału~5 daje ciągłą funkcję
$f:X\to[0,1]$ równą $1$ na $F$
i $0$ na $X\setminus W$.
Zatem $\{f\ne0\}\subseteq W$,
a po domknięciu
$\operatorname{supp}f
\subseteq\overline W^{\,X}\subseteq V$.
\end{proof}

\section{Konstrukcja przeliczalnego podziału jedności}
\label{sec:topog-parazwartosc-konstrukcja}

\begin{twierdzenie}[Podział jedności w regularnej przestrzeni Lindelöfa]
\label{thm:topog-parazwartosc-lindelof}
Niech $X$ będzie regularną przestrzenią
Hausdorffa o własności Lindelöfa,
a $\mathcal U$ dowolnym pokryciem
otwartym $X$. Istnieje co najwyżej
przeliczalny, lokalnie skończony
podział jedności na $X$
podporządkowany $\mathcal U$.
W konsekwencji $X$ jest parazwarta.
\end{twierdzenie}
\begin{proof}
Dla każdego $x\in X$ wybieramy
$U_x\in\mathcal U$ z $x\in U_x$.
Z regularności i kryterium zwężania
otoczeń z rozdziału~5 otrzymujemy
otwarte zbiory $W_x,V_x$ takie, że
\[
x\in W_x\subseteq\overline{W_x}^{\,X}
\subseteq V_x\subseteq
\overline{V_x}^{\,X}\subseteq U_x.
\]
Rodzina $(W_x)_{x\in X}$ pokrywa $X$.
Z własności Lindelöfa wybieramy
przeliczalne podpokrycie
$(W_n)_{n\ge1}$ z odpowiadającymi
zbiorami $V_n$ i $U_n$.
Jeśli wybrane podpokrycie jest skończone,
dopowiadamy puste $W_n,V_n$ i dowolne
$U_n\in\mathcal U$ dla pozostałych
indeksów; przy $X=\varnothing$ twierdzenie
jest bezpośrednie.

Z rozdziału~10 wiemy, że przestrzeń
regularna i Lindelöfa jest normalna.
W związku z tym
lemat~\ref{lem:topog-parazwartosc-odciecia}
zastosowany do
$F_n=\overline{W_n}^{\,X}$,
otwartego zbioru $U_n$ i otwartego
zbioru pośredniego $V_n$ daje funkcje ciągłe
$f_n:X\to[0,1]$ takie, że
$f_n|_{F_n}=1$ i
$\operatorname{supp}f_n\subseteq
\overline{V_n}^{\,X}\subseteq U_n$.
Definiujemy kolejno
\[
\varphi_1=f_1,\qquad
\varphi_n=f_n\prod_{j=1}^{n-1}(1-f_j),
\quad n\ge2.
\]
Każda $\varphi_n$ jest ciągła,
nieujemna i nie przekracza $1$.
Ponadto $\{x:\varphi_n(x)\ne0\}
\subseteq\{x:f_n(x)\ne0\}$,
więc
$\operatorname{supp}\varphi_n
\subseteq\operatorname{supp}f_n
\subseteq U_n$.

Sprawdzimy lokalną skończoność nośników,
a nie tylko punktową skończoność sumy.
Ustalmy $x\in X$.
Ponieważ $(W_n)$ pokrywa $X$,
istnieje $N$ z $x\in W_N$.
Na całym otwartym otoczeniu $W_N$
funkcja $f_N$ jest równa $1$.
Dla $n>N$ we wzorze na $\varphi_n$
występuje czynnik $1-f_N$,
więc $\varphi_n$ zeruje się na $W_N$.
Wobec otwartości $W_N$ również
$\operatorname{supp}\varphi_n
\cap W_N=\varnothing$ dla $n>N$.
Otoczenie $W_N$ spotyka zatem
co najwyżej $N$ nośników.

Dla każdego $m\ge1$ zachodzi
tożsamość teleskopowa
\[
\sum_{n=1}^{m}\varphi_n(x)
=1-\prod_{n=1}^{m}(1-f_n(x)).
\]
Można ją sprawdzić przez indukcję,
odejmując kolejne iloczyny.
Jeżeli $x\in W_N$, to $f_N(x)=1$,
a więc dla wszystkich $m\ge N$
prawa strona jest równa $1$.
To dowodzi warunku sumy.
Otrzymaliśmy żądany podział jedności;
parazwartość wynika ze
stwierdzenia~\ref{prop:topog-parazwartosc-jednosc-pokrycie}.
\end{proof}

\begin{wniosek}[Dwie klasy przestrzeni parazwartych]
\label{cor:topog-parazwartosc-klasy}
Każda przestrzeń metryzowalna
spełniająca drugi aksjomat przeliczalności
oraz każda lokalnie zwarta, $\sigma$-zwarta
przestrzeń Hausdorffa jest parazwarta
i ma podział jedności podporządkowany
dowolnemu pokryciu otwartemu.
\end{wniosek}
\begin{proof}
W pierwszym przypadku przestrzeń jest
regularna i Hausdorffa jako metryzowalna,
a z rozdziału~9 spełnia własność Lindelöfa.
W drugim przypadku lokalna zwartość
i Hausdorffowość zapewniają regularność:
lemat o otoczeniach mających zwarte
domknięcia z rozdziału~14 daje
dla dowolnego otwartego otoczenia
punktu mniejsze otwarte otoczenie
o domknięciu zawartym w wyjściowym.
$\sigma$-zwartość implikuje własność
Lindelöfa: każde pokrycie otwarte
ma skończone podpokrycie każdego
ze zwartych zbiorów tworzących
przeliczalną sumę $X$; suma tych
skończonych rodzin jest przeliczalna.
W obu przypadkach stosujemy
twierdzenie~\ref{thm:topog-parazwartosc-lindelof}.
\end{proof}

\begin{przyklad}[Jawny podział na prostej]
\label{ex:topog-parazwartosc-prosta}
Dla $n\in\mathbb Z$ i $x\in\R$ połóżmy
$h_n(x)=\max\{0,1-|x-n|\}$.
Są to funkcje ciągłe o nośnikach
$[n-1,n+1]$. Rodzina nośników
jest lokalnie skończona.
Jeśli $x\in[n,n+1]$, to
$h_n(x)=n+1-x$ i
$h_{n+1}(x)=x-n$,
a wszystkie pozostałe $h_k(x)$
są równe zeru. Stąd
$\sum_{n\in\mathbb Z}h_n(x)=1$.
Otrzymujemy podział jedności
podporządkowany pokryciu
$((n-\tfrac32,n+\tfrac32))_{n\in\mathbb Z}$.
Warto zauważyć, że w punktach całkowitych
jeden składnik ma wartość $1$,
a wszystkie pozostałe wartość $0$.
\end{przyklad}

\begin{przyklad}[Nieprzeliczalna przestrzeń dyskretna]
\label{ex:topog-parazwartosc-dyskretna}
Niech $D$ będzie dowolną, także
nieprzeliczalną przestrzenią dyskretną.
Jest ona parazwarta: z danego
pokrycia otwartego wybieramy dla każdego
$x\in D$ singleton $\{x\}$.
Rodzina singletonów jest lokalnie skończona.
Funkcje $\chi_{\{x\}}:D\to[0,1]$
są ciągłe i tworzą podporządkowany
temu pokryciu podział jedności.
Przykład pokazuje, że parazwartość
nie implikuje własności Lindelöfa.
\end{przyklad}

\section{Twierdzenie Stone'a}
\label{sec:topog-parazwartosc-stone}

\begin{twierdzenie}[Twierdzenie Stone'a]
\label{thm:topog-parazwartosc-stone}
Niech $(X,d)$ będzie przestrzenią metryczną.
Wówczas $X$ jest parazwarta.
\end{twierdzenie}
\begin{proof}
Jeżeli $X=\varnothing$, to pusta rodzina
jest żądanym uszczegółowieniem.
Niech zatem $X\ne\varnothing$.
Korzystając z twierdzenia o dobrym uporządkowaniu,
zapiszmy pokrycie jako
$\mathcal U=(U_\alpha)_{\alpha<\kappa}$,
gdzie $\kappa$ jest liczbą porządkową.
Dla $x\in X$ istnieje najmniejszy
indeks $\alpha(x)<\kappa$ taki, że
$x\in U_{\alpha(x)}$.
Piszemy $r_n=2^{-n}$ dla $n\ge1$.

Konstruujemy kolejno, dla $n=1,2,\ldots$,
otwarte zbiory $V_{\alpha,n}$,
$\alpha<\kappa$.
Po zdefiniowaniu wszystkich poziomów
o indeksach mniejszych od $n$ określamy
zbiór dopuszczalnych środków
\[
E_{\alpha,n}=
\left\{z\in X:
\begin{array}{l}
\alpha(z)=\alpha,\\
z\notin\displaystyle\bigcup_{j<n}
            \bigcup_{\beta<\kappa}V_{\beta,j},\\
B_d(z,3r_n)\subseteq U_\alpha
\end{array}\right\}
\]
i przyjmujemy
\[
V_{\alpha,n}=\bigcup_{z\in E_{\alpha,n}}B_d(z,r_n).
\]
Dla $n=1$ suma po $j<n$ jest pusta.
Każdy $V_{\alpha,n}$ jest otwarty,
a z trzeciego warunku
$V_{\alpha,n}\subseteq U_\alpha$.

\emph{Pokrywanie.}
Niech $x\in X$ i $\alpha=\alpha(x)$.
Ponieważ $U_\alpha$ jest otwarty,
istnieje $n\ge1$ z
$B_d(x,3r_n)\subseteq U_\alpha$.
Jeżeli $x$ należy już do któregoś
$V_{\beta,j}$ z $j<n$, to jest pokryty.
W przeciwnym razie $x\in E_{\alpha,n}$,
a stąd $x\in V_{\alpha,n}$.
Zatem rodzina wszystkich
$V_{\alpha,n}$ pokrywa $X$.

\emph{Rozdzielenie na ustalonym poziomie.}
Ustalmy $i\ge1$ i dwa różne indeksy
$\beta<\gamma<\kappa$.
Jeżeli
$p\in V_{\beta,i}$ i
$q\in V_{\gamma,i}$,
to istnieją środki
$y\in E_{\beta,i}$,
$z\in E_{\gamma,i}$
z $d(p,y)<r_i$ i $d(q,z)<r_i$.
Ponieważ $\alpha(z)=\gamma>\beta$,
punkt $z$ nie należy do $U_\beta$.
Jednocześnie
$B_d(y,3r_i)\subseteq U_\beta$,
zatem $d(y,z)\ge3r_i$.
Nierówność trójkąta daje
\[
d(p,q)\ge d(y,z)-d(p,y)-d(q,z)>r_i.
\]
W konsekwencji otwarta kula
o średnicy nie większej niż $r_i$
może spotkać najwyżej jeden
z różnych zbiorów $V_{\beta,i}$.
Istotnie, gdyby spotykała dwa,
odległość wybranych w nich punktów
byłaby mniejsza od jej średnicy,
co przeczyłoby powyższemu oszacowaniu.

\emph{Lokalna skończoność wszystkich poziomów.}
Ustalmy $x\in X$. Z wykazanego
pokrywania wybierzmy $\alpha$ i $n$
takie, że $x\in V_{\alpha,n}$.
Z otwartości tego zbioru istnieje
$\rho>0$ z
$B_d(x,\rho)\subseteq V_{\alpha,n}$.
Wybierzmy liczbę całkowitą $M>n$
tak dużą, że $2r_M<\rho$,
i połóżmy $O=B_d(x,r_M)$.

Jeśli $i\ge M$ i
$z\in E_{\beta,i}$, to
drugi warunek definicji środka
wyklucza $z\in V_{\alpha,n}$,
ponieważ $n<i$.
W szczególności $d(x,z)\ge\rho$.
Dla $p\in B_d(z,r_i)$ mamy
\[
d(x,p)\ge d(x,z)-d(z,p)
>\rho-r_i\ge\rho-r_M>r_M.
\]
Zatem $O\cap V_{\beta,i}=\varnothing$
dla wszystkich $i\ge M$ i
$\beta<\kappa$.

Dla każdego z pozostałych poziomów
$1\le i<M$ zachodzi
$r_i\ge 2r_M$, a więc
średnica kuli $O$ jest nie większa
niż $r_i$. Na mocy wykazanego
rozdzielenia $O$ spotyka
co najwyżej jeden zbiór
$V_{\beta,i}$ na każdym takim poziomie.
W sumie $O$ spotyka najwyżej
$M-1$ zbiorów indeksowanej
rodziny $(V_{\beta,i})_{\beta<\kappa,\ i\ge1}$.
Dowodzi to lokalnej skończoności.

Otrzymana rodzina jest otwartym
pokryciem, jest lokalnie skończona
i uszczegóławia $\mathcal U$.
Ponieważ pokrycie $\mathcal U$ było
dowolne, przestrzeń $X$ jest parazwarta.
\end{proof}

\begin{uwaga}[Znaczenie założeń]
Dowód nie wymaga zupełności, ośrodkowości,
zwartości lokalnej ani żadnego aksjomatu
przeliczalności. Wykorzystuje
dobry porządek indeksów pokrycia,
zatem w standardowym ujęciu odwołuje
się do aksjomatu wyboru. W odróżnieniu
od wcześniejszej konstrukcji dla
przestrzeni Lindelöfa uzyskano tutaj
od razu lokalnie skończone pokrycie
otwarte, a nie podział jedności.
\end{uwaga}


\section{Twierdzenie Nagaty--Smirnowa o metryzowalności}
\label{sec:topog-parazwartosc-nagata-smirnow}

\begin{definicja}[Baza będąca przeliczalną sumą rodzin lokalnie skończonych]
\label{def:topog-parazwartosc-sigma-lokalnie-skonczona-baza}
Niech $X$ będzie przestrzenią topologiczną.
Rodzinę $\mathcal B$ jej otwartych podzbiorów
nazywamy \emph{bazą będącą przeliczalną sumą
rodzin lokalnie skończonych}, jeżeli jest bazą
topologii i istnieją rodziny
$\mathcal B_n$, $n\ge1$, lokalnie skończone
w $X$, takie, że
$\mathcal B=\bigcup_{n\ge1}\mathcal B_n$.
Nie zakładamy, że poszczególne
$\mathcal B_n$ są przeliczalne.
\end{definicja}

\begin{twierdzenie}[Twierdzenie Nagaty--Smirnowa]
\label{thm:topog-parazwartosc-nagata-smirnow}
Niech $X$ będzie przestrzenią topologiczną.
Następujące warunki są równoważne:
\begin{enumerate}[label=(\roman*)]
\item $X$ jest metryzowalna;
\item $X$ jest regularną przestrzenią $T_1$
      i ma bazę będącą przeliczalną sumą
      rodzin lokalnie skończonych.
\end{enumerate}
\end{twierdzenie}
\begin{proof}
\emph{Konieczność.}
Załóżmy, że topologię $X$ wyznacza
metryka $d$. Przestrzeń metryczna
jest regularna i spełnia aksjomat $T_1$
(rozdział~10). Dla każdego $n\ge1$
rozpatrzmy pokrycie otwarte
$\mathcal U_n=\{B_d(x,2^{-n}):x\in X\}$.
Na mocy twierdzenia Stone'a,
udowodnionego powyżej, istnieje
lokalnie skończone otwarte pokrycie
$\mathcal V_n$ uszczegóławiające
$\mathcal U_n$.

Pokażemy, że
$\mathcal B=\bigcup_{n\ge1}\mathcal V_n$
jest bazą topologii.
Niech $x\in O$, gdzie $O$ jest otwarty.
Wybierzmy $\varepsilon>0$
z $B_d(x,\varepsilon)\subseteq O$
oraz $n$ z $2^{1-n}<\varepsilon$.
Ponieważ $\mathcal V_n$ pokrywa $X$,
istnieje $V\in\mathcal V_n$ z $x\in V$.
Z uszczegółowienia otrzymujemy
$y\in X$ takie, że
$V\subseteq B_d(y,2^{-n})$.
Wówczas $d(x,y)<2^{-n}$,
więc dla każdego $z\in V$
mamy $d(x,z)<2^{1-n}<\varepsilon$.
Zatem $x\in V\subseteq O$.
Każda $\mathcal V_n$ jest lokalnie
skończona, co dowodzi (ii).

\emph{Wystarczalność.}
Przyjmijmy (ii) i zapiszmy
$\mathcal B=\bigcup_{n\ge1}\mathcal B_n$,
gdzie każda $\mathcal B_n$ jest lokalnie
skończona. Najpierw przygotujemy
dwa fakty potrzebne do zbudowania metryki.

\emph{Krok 1: normalność.}
Niech $A,C\subseteq X$ będą rozłączne
i domknięte. Dla każdego $n$ określmy
\[
\mathcal P_n=
\{B\in\mathcal B_n:
\overline B^{\,X}\cap C=\varnothing\},
\qquad
P_n=\bigcup_{B\in\mathcal P_n}B,
\]
a analogicznie
\[
\mathcal Q_n=
\{B\in\mathcal B_n:
\overline B^{\,X}\cap A=\varnothing\},
\qquad
Q_n=\bigcup_{B\in\mathcal Q_n}B.
\]
Każdy $P_n,Q_n$ jest otwarty.
Lokalna skończoność $\mathcal B_n$
oraz stwierdzenie
\ref{prop:topog-parazwartosc-domkniecia}
dają
$\overline{P_n}^{\,X}
=\bigcup_{B\in\mathcal P_n}\overline B^{\,X}$,
stąd
$\overline{P_n}^{\,X}\cap C=\varnothing$.
Podobnie
$\overline{Q_n}^{\,X}\cap A=\varnothing$.
Z regularności i własności bazy
wynika, że $(P_n)$ pokrywa $A$,
a $(Q_n)$ pokrywa $C$:
dla $a\in A$ można wybrać
otwarty $H\ni a$ z
$\overline H^{\,X}\subseteq X\setminus C$,
a następnie $B\in\mathcal B$
z $a\in B\subseteq H$;
wtedy $B\in\mathcal P_n$ dla pewnego $n$.
Dla punktów $C$ postępujemy analogicznie.

Połóżmy
\[
P=\bigcup_{n\ge1}
 \left(P_n\setminus
       \bigcup_{j\le n}\overline{Q_j}^{\,X}\right),
\qquad
Q=\bigcup_{n\ge1}
 \left(Q_n\setminus
       \bigcup_{j\le n}\overline{P_j}^{\,X}\right).
\]
Oba zbiory są otwarte, gdyż każdorazowo
odejmujemy skończoną sumę zbiorów
domkniętych od otwartego.
Mamy $A\subseteq P$, ponieważ
$A\cap\overline{Q_j}^{\,X}=\varnothing$
dla każdego $j$, i podobnie
$C\subseteq Q$.
Zbiory $P$ i $Q$ są rozłączne:
gdyby $x$ należał do składnika
$P_n$ oraz składnika $Q_m$,
to dla $m\le n$ mielibyśmy
$x\notin\overline{Q_m}^{\,X}$,
wbrew $x\in Q_m$, a dla $n\le m$
otrzymalibyśmy analogiczną
sprzeczność z $x\in P_n$.
Zatem $X$ jest normalna.

\emph{Krok 2: funkcje dodatnie dokładnie
na zadanym zbiorze otwartym.}
Niech $G\subseteq X$ będzie otwarty.
Dla każdego $n$ określmy
\[
F_n(G)=
\bigcup\{\overline B^{\,X}:
 B\in\mathcal B_n,\
 \overline B^{\,X}\subseteq G\}.
\]
Jest to zbiór domknięty:
rodzina występujących tu $B$
jest lokalnie skończona, więc
jej domknięcia mają domkniętą sumę,
zgodnie z wcześniejszym stwierdzeniem.
Ponadto $F_n(G)\subseteq G$.
Z regularności i własności bazy
wynika
$G=\bigcup_{n\ge1}F_n(G)$,
gdyż dla $x\in G$ można wybrać
$B\in\mathcal B$ z
$x\in B\subseteq\overline B^{\,X}
\subseteq G$.

Z kroku 1 i lematu Urysohna
(rozdział~5) dla każdego $n$
istnieje funkcja ciągła
$g_n:X\to[0,1]$ taka, że
$g_n=1$ na $F_n(G)$ i
$g_n=0$ na $X\setminus G$.
Szereg
\[
g_G(x)=\sum_{n\ge1}2^{-n}g_n(x)
\]
jest jednostajnie zbieżny,
ponieważ $0\le g_n\le1$ i
$\sum_{n\ge1}2^{-n}=1$.
Wynika stąd ciągłość $g_G$.
Dla $x\notin G$ wszystkie składniki
są zerowe; dla $x\in G$
co najmniej jeden składnik jest
równy $2^{-n}$.
Zatem $g_G(x)>0$ dokładnie dla $x\in G$.

\emph{Krok 3: konstrukcja metryki.}
Dla każdego $n\ge1$ i
$B\in\mathcal B_n$ wybierzmy
funkcję $g_B$ z kroku 2 i połóżmy
$f_{n,B}=2^{-n}g_B$.
Wtedy $f_{n,B}:X\to[0,2^{-n}]$
jest ciągła oraz
\[
f_{n,B}(x)>0
\ \Longleftrightarrow\ x\in B.
\]
Niech
$J=\{(n,B):n\ge1,\ B\in\mathcal B_n\}$.
Zbiór indeksów $J$ nie musi być
przeliczalny. Definiujemy
$F:X\to\ell^\infty(J)$ wzorem
$F(x)=(f_{n,B}(x))_{(n,B)\in J}$.
Jest to poprawnie określone, ponieważ
$0\le f_{n,B}(x)\le1$.

Pokażemy ciągłość $F$ względem
normy supremum. Ustalmy $x\in X$
i $\varepsilon>0$.
Wybierzmy $N$ z
$2^{-N}<\varepsilon/2$.
Dla $n\ge N$ i dowolnego $B$
różnica
$|f_{n,B}(x)-f_{n,B}(y)|$
jest mniejsza niż $\varepsilon/2$
dla każdego $y\in X$.
Dla każdego z skończenie wielu $n<N$
lokalna skończoność $\mathcal B_n$
daje otwarte otoczenie $O_n\ni x$
przecinające jedynie skończenie
wiele zbiorów z $\mathcal B_n$.
Poza tymi zbiorami wszystkie
$f_{n,B}$ zerują się na $O_n$
oraz w punkcie $x$.
Dla pozostałych, skończenie
wielu funkcji ciągłość w $x$
pozwala dodatkowo zwęzić otoczenie,
tak by
$|f_{n,B}(y)-f_{n,B}(x)|
<\varepsilon/2$.
Przecięcie skończenie wielu
takich otoczeń dla $n<N$
jest otwartym otoczeniem $O$ punktu $x$,
a dla $y\in O$ zachodzi
\[
\|F(y)-F(x)\|_\infty<\varepsilon.
\]
Stąd $F$ jest ciągła.

Ponieważ $X$ spełnia aksjomat $T_1$,
dla różnych $x,y\in X$ istnieje
otwarty zbiór $G$ taki, że
$x\in G$, $y\notin G$.
Wybierzmy $B\in\mathcal B_n$
z $x\in B\subseteq G$.
Wtedy $f_{n,B}(x)>0$ i
$f_{n,B}(y)=0$, zatem $F(x)\ne F(y)$.

Pozostaje sprawdzić ciągłość
odwzorowania odwrotnego
$F^{-1}:F(X)\to X$.
Niech $G\subseteq X$ będzie otwarty
i $x\in G$. Wybierzmy
$B\in\mathcal B_n$ z
$x\in B\subseteq G$ i połóżmy
$\delta=f_{n,B}(x)/2>0$.
Jeśli $y\in X$ i
$\|F(y)-F(x)\|_\infty<\delta$,
to $f_{n,B}(y)>0$, a więc
$y\in B\subseteq G$.
To dowodzi, że obraz
$F(G)$ jest otwarty w $F(X)$.
Zatem $F$ jest homeomorfizmem
$X$ na podprzestrzeń metrycznej
przestrzeni $\ell^\infty(J)$.
Metrykę zgodną na $X$ otrzymujemy
jako
\[
d(x,y)=
\sup_{(n,B)\in J}
|f_{n,B}(x)-f_{n,B}(y)|.
\]
To kończy dowód równoważności.
\end{proof}

\begin{uwaga}[Porównanie z kryterium Urysohna]
\label{rem:topog-parazwartosc-nagata-urysohn}
Twierdzenie Urysohna z rozdziału~10
wymaga regularności, aksjomatu $T_1$
i drugiego aksjomatu przeliczalności.
Baza przeliczalna jest automatycznie
przeliczalną sumą rodzin
jednoelementowych, zatem twierdzenie
Nagaty--Smirnowa obejmuje ten wynik.
Odwrotna implikacja nie zachodzi:
nieprzeliczalna przestrzeń dyskretna
ma bazę singletonów, która sama jest
lokalnie skończona, ale nie spełnia
drugiego aksjomatu przeliczalności.
Główną różnicą w dowodzie jest
użycie przestrzeni $\ell^\infty(J)$
dla dowolnego, także nieprzeliczalnego
zbioru indeksów $J$, z wagami
$2^{-n}$ kontrolującymi ciągłość.
\end{uwaga}
